% Síntesis original a partir de la guía docente ULE 2026-2027,
% asignatura 0710311; no es una publicación oficial de la ULE.
\section*{Cómo utilizar este manual según la guía docente 2026--2027}
\label{sec:orientacion-ule-2026}
La guía firmada de la Universidad de León para \emph{Mecánica de Fluidos}
(0710311; segundo curso, primer semestre; 6 ECTS) distribuye la materia
en seis bloques. La correspondencia entre los bloques de la guía y
los capítulos de este libro es la siguiente. Se ofrece como itinerario
de estudio, no como certificación de cada afirmación del manuscrito.

\begin{enumerate}
\item \textbf{Tema 1: fundamentos de mecánica clásica.} Capítulos 1 y 2:
leyes de Newton, momentos, referencias no inerciales, energía,
Lagrange y sistemas caóticos.
\item \textbf{Tema 2: fundamentos de dinámica de fluidos.} Capítulos 3 y 4:
medio continuo, esfuerzos, descripciones lagrangiana y euleriana,
derivada material y condiciones de contorno.
\item \textbf{Tema 3: estudio de casos.} Capítulos 5 al 9:
hidrostática, Couette y Poiseuille, fuerzas sobre cuerpos,
oscilaciones, motores a reacción y aerogeneradores.
\item \textbf{Tema 4: similaridad dinámica.} Capítulo 10:
semejanza, números adimensionales, Reynolds, fricción de
Darcy--Weisbach y teorema Pi.
\item \textbf{Tema 5: formalismo matemático y aplicaciones.} Capítulos
11 al 13: operadores, teoremas integrales, transporte de Reynolds,
continuidad, Bernoulli, difusión y tensores de esfuerzo y deformación.
\item \textbf{Tema 6: Navier--Stokes.} Capítulos 14 al 16:
ecuaciones de conservación, Euler y Navier--Stokes,
adimensionalización, Stokes, Boussinesq y sistemas giratorios.
\end{enumerate}

\noindent\textbf{Carga de trabajo declarada por la guía.}
Son 150 horas en total: 60 presenciales y 90 de trabajo autónomo.
Las presenciales se dividen en 27 horas de clases magistrales,
12 de resolución de ejercicios y 21 de prácticas informáticas.
Este desglose corresponde a la planificación docente, no a una
distribución proporcional de páginas ni a una ponderación del examen.

\noindent\textbf{Preparación de la evaluación.}
El examen escrito final aporta el 70\,\% y los ejercicios realizados
\emph{presencialmente durante el curso} el 30\,\%.
La guía exige obtener al menos \(5/10\) en el examen para aprobar.
Si el examen queda por debajo de 5, la calificación de la convocatoria
es la de ese examen; con un examen de 5 o más, la calificación se
calcula con la ponderación anterior, garantizando como mínimo 5.
Los ejercicios de clase cuentan también en segunda convocatoria del
mismo curso, pero no en cursos posteriores. Por eso es útil estudiar
cada capítulo en tres pasadas: comprender y deducir el modelo;
resolver problemas sin consultar la solución; y reproducir,
interpretar y criticar un cálculo mediante el ordenador.

\noindent\textbf{Sobre los materiales de apoyo.}
Las siete prácticas Mathematica, los 48 problemas del cuaderno
complementario y los cuatro simulacros son recursos
\emph{de este proyecto}, no materiales oficiales de evaluación.
En particular, la guía no fija una duración ni una cantidad de
problemas en el examen real, y el profesor comunica antes de las pruebas
los materiales y recursos que se permiten.

\noindent\textit{Fuente:} Universidad de León, \emph{Guía Docente
2026--2027, 0710311, Mecánica de Fluidos}, páginas 2--5; documento
firmado electrónicamente el 18 de agosto de 2026.
La trazabilidad detallada de la correspondencia temática se conserva
en \texttt{bibliography/guia-docente-ule-2026-2027.json}.
